This project addresses the fundamental challenge of extending the operational endurance of autonomous marine vehicles. We propose a battery-aware reinforcement learning policy that integrates sensor management with motion planning, explicitly adapting to environmental dynamics like currents and waves. By treating energy conservation as a first-class objective alongside exploration, we aim to demonstrate that intelligent, context-aware behaviors can significantly increase the effective range and utility of ASVs.

We benchmark this approach against standard frontier and rule-based methods, as well as non-adaptive RL baselines. The expected outcome is a policy that not only covers more area per Joule but also exhibits distinct, interpretable strategies for different water bodies—such as exploiting currents for transport or duty-cycling sensors in safe regions. If full RL training proves unstable, a fallback to a hierarchical approach with separate planning and power management modules will still provide valuable insights into the energy dynamics of marine autonomy.
